这些练习要求把模型条件转成可核对的等式与不等式。观测矩题先检查同一潜在模型能否同时满足全部协方差关系；极值概率题则把“证明上界”与“达到上界”分开处理，严格事件的边界质量尤其重要。

约束估计先列可行域，再检查乘子符号、互补松弛与唯一性。预算把坐标联系起来，不能只做各自的阈值运算。得到候选解后，可用直接代入、单变量消元或目标值增量作独立核对，最后说明结论适用的参数范围。
\par\smallskip
